\documentclass[10pt,a4paper]{amsart}
\usepackage[margin=19mm]{geometry}
\usepackage{amsmath,amssymb,mathtools}
\usepackage[scr=rsfs,cal=euler]{mathalfa}
\usepackage{xfrac}
\usepackage[unicode,hidelinks,bookmarks=false]{hyperref}
\newcommand{\ie}{i.e.\ }
\newcommand{\cf}{cf.\ }
\newcommand{\resp}{respectively\ }
\newcommand{\Subsection}{\subsection}
\providecommand{\nobreakdash}{\nobreak\hbox{-}}
\setlength{\emergencystretch}{2em}
\multlinegap0pt
\usepackage{bm}
\usepackage{bbm}
\usepackage{dsfont}
\usepackage{xspace}
\usepackage{cancel}
\usepackage{mathtools}
\usepackage{amsthm}
\makeatletter
\@ifundefined{theorem}{\newtheorem{theorem}{Theorem}}{}
\@ifundefined{prop}{\newtheorem{prop}[theorem]{Proposition}}{}
\makeatother
\title[A Rigorous Functional-Integral Construction of Toral Chern-S]{A Rigorous Functional-Integral Construction of Toral Chern-Simons Theory}
\author{Daniel Galviz}
\date{Source: arXiv:2604.02013v1 (CC BY 4.0)}
\begin{document}
\maketitle

\noindent\emph{Source and licence.} A Rigorous Functional-Integral Construction of Toral Chern-Simons Theory. Version arXiv:2604.02013v1 (CC BY 4.0), distributed under the Creative Commons Attribution 4.0 International licence, as recorded in the arXiv licence metadata for this exact version and shown on the abstract page https://arxiv.org/abs/2604.02013v1. Full rights notice: licenses/arxiv-2604.02013-CC-BY-4.0.txt.\par\smallskip

\noindent\textit{Editorial scope.} Counted unit: the Prop whose source label is \texttt{prop:closed-gauge-reduction} in the article (source0.tex lines 710--729, source label \texttt{prop:closed-gauge-reduction}), delivered together with the article's own complete original proof of it (source0.tex lines 731--849, one \texttt{proof} environment of 119 lines). No mathematics is written, repaired or re-derived: statement and proof are the article's own text line for line. Required context retained uncounted: prop:moduli-torsor-toral (lines 443--449); prop:linear-quotient-factorization (lines 629--649). Every definition, display and label of those lines is retained unchanged. Omitted from this package and not redistributed: 1--442, 450--628, 650--709, 730--730, 850--1806. Nothing inside a retained excerpt is omitted or altered: every retained body is delivered exactly as the article's lines are written, so no omission receipt of any kind accompanies this package. Citations: the retained excerpts carry no citation command at all, so no bibliography entry of the article is transcribed and no reference is redistributed. Reference adaptation: none was needed and none was made; every reference inside the delivered unit resolves inside it, so no unresolved reference or citation is produced. Printed slips kept literal: no printed word, symbol, hypothesis or number of the counted statement or of its proof was altered. Layout and class: the article's own class is \texttt{amsart} with packages including geometry, fontenc, babel, microtype, tikz, tikz-cd, amsmath, amssymb, amsthm, amscd, mathtools, mathrsfs, enumitem, xcolor; the wrapper is the lane's pinned \texttt{amsart} editorial wrapper at 10pt on A4, which restates the theorem-like environments the retained text needs and adds the article's own macro definitions that the retained text needs (none), with extra wrapper packages (bm, bbm, dsfont, xspace, cancel, mathtools). The delivered numbering is the unit's own. Assets: no retained span contains a figure environment, a graphics-inclusion command, a PDF-page inclusion, a table or a TikZ picture; no image file is copied into this package, the package declares no asset dependency and no image file is redistributed with it. Licence: version arXiv:2604.02013v1 of this article is distributed under the Creative Commons Attribution 4.0 International licence, as recorded in the arXiv licence metadata for this exact version and shown on the abstract page https://arxiv.org/abs/2604.02013v1. A full rights notice is delivered at licenses/arxiv-2604.02013-CC-BY-4.0.txt. Characters: every retained excerpt is pure ASCII, so the delivered engine, XeLaTeX with its default Latin Modern text font, reports no missing character.
\begin{prop}
\label{prop:moduli-torsor-toral}
(i)If \(X\) is closed, then  \(\mathcal M_{X,p}(\mathbb T)\) is a torsor for $H^1(X;\mathfrak t)/H^1(X;\Lambda).$

\noindent(ii) If \(\partial X=\Sigma\neq\varnothing\) and \(\eta\) is extendable, then \(\mathcal M_{X,p}(\eta;\mathbb T)\) is a torsor for \\
$H^1(X,\Sigma;\mathfrak t)/H^1(X,\Sigma;\Lambda).$
\end{prop}

\begin{prop}
\label{prop:linear-quotient-factorization}
Let \(V\) be a finite-dimensional Euclidean space and suppose that
\[
V=H\oplus E\oplus N
\]
is an orthogonal decomposition. Let \(\Gamma\subset H\) be a full lattice, and let
\(A\colon N\to N\) be self-adjoint and invertible. Define a function \(q\colon V\to \mathbb C\) by
\[
q(h+e+n)=q_0+\frac12\langle An,n\rangle,
\qquad
h\in H,\ e\in E,\ n\in N,
\]
for some constant \(q_0\in \mathbb C\). Then the oscillatory integral on the quotient by
\(E\), with periodicity lattice \(\Gamma\) in the \(H\)-direction, factors as
\[
\int_{(H/\Gamma)\times N} e^{iq(h+n)}\,d(h+\Gamma)\,dn
=
e^{iq_0}\,\operatorname{vol}(H/\Gamma)\,\operatorname{Fres}(A).
\]
\end{prop}

\begin{prop}
\label{prop:closed-gauge-reduction}
With the normalization conventions of  rank-one functional-integral construction,
the exact Gaussian evaluation of the formal quotient integral in the torsion class \(p\) is
\begin{equation}
\label{eq:closed-gauge-reduction}
\mathcal Z^{\mathrm{gauss}}_{X,p}(\mathbb T,K)
=
e^{2\pi i\,\operatorname{CS}_{X,P,K}(\Theta_p)}\,
\det\nolimits'_{\zeta}(\Delta_{0,\mathfrak t})^{1/2}\,
\mathcal I_\zeta(B_{X,K})\,
\int_{\mathcal M_{X,p}(\mathbb T)} d\mu_{\mathrm{har}},
\end{equation}
where \(d\mu_{\mathrm{har}}\) denotes the translation-invariant density on the harmonic torus
\[
\mathcal M_{X,p}(\mathbb T)
\cong
H^1(X;\mathfrak t)/H^1(X;\Lambda).
\]
\end{prop}

\begin{proof}
Choose a Riemannian metric on \(X\). By Hodge decomposition,
\[
\Omega^1(X;\mathfrak t)
=
H^1(X;\mathfrak t)
\oplus
d\Omega^0(X;\mathfrak t)
\oplus
d^*\Omega^2(X;\mathfrak t).
\]
Accordingly, every \(a\in \Omega^1(X;\mathfrak t)\) can be written uniquely as
\[
a=a_{\mathrm{har}}+d\phi+a_{\mathrm{coex}},
\qquad
a_{\mathrm{har}}\in H^1(X;\mathfrak t),\quad
\phi\in \Omega^0(X;\mathfrak t),\quad
a_{\mathrm{coex}}\in d^*\Omega^2(X;\mathfrak t).
\]

Let
\[
q(a):=\frac{i}{4\pi}\int_X K(a\wedge da).
\]
Since \(da_{\mathrm{har}}=0\) and \(d^2\phi=0\), we have $da=da_{\mathrm{coex}}.$ Expanding \(q(a)\) gives
\[
q(a)
=
\frac{i}{4\pi}\int_X
\Bigl(
K(a_{\mathrm{har}}\wedge da_{\mathrm{coex}})
+
K(d\phi\wedge da_{\mathrm{coex}})
+
K(a_{\mathrm{coex}}\wedge da_{\mathrm{coex}})
\Bigr).
\]
The first two terms vanish. Indeed, since \(X\) is closed,
\[
d\,K(a_{\mathrm{har}}\wedge a_{\mathrm{coex}})
=
K(da_{\mathrm{har}}\wedge a_{\mathrm{coex}})
-
K(a_{\mathrm{har}}\wedge da_{\mathrm{coex}})
=
-
K(a_{\mathrm{har}}\wedge da_{\mathrm{coex}}),
\]
so
\[
\int_X K(a_{\mathrm{har}}\wedge da_{\mathrm{coex}})=0.
\]
Similarly,
\[
d\,K(d\phi\wedge a_{\mathrm{coex}})
=
K(d^2\phi\wedge a_{\mathrm{coex}})
-
K(d\phi\wedge da_{\mathrm{coex}})
=
-
K(d\phi\wedge da_{\mathrm{coex}}),
\]
hence
\[
\int_X K(d\phi\wedge da_{\mathrm{coex}})=0.
\]
Therefore
\[
q(a)=\frac{i}{4\pi}\int_X K(a_{\mathrm{coex}}\wedge da_{\mathrm{coex}}).
\]

Equivalently, if \(B_{X,K}\) denotes the self-adjoint operator on
\(d^*\Omega^2(X;\mathfrak t)\) defined earlier, then
\[
q(a)=\frac{i}{4\pi}\,\langle B_{X,K}a_{\mathrm{coex}},a_{\mathrm{coex}}\rangle.
\]
Moreover, \(B_{X,K}\) is nondegenerate on \(d^*\Omega^2(X;\mathfrak t)\): if
\(B_{X,K}a_{\mathrm{coex}}=0\), then \(da_{\mathrm{coex}}=0\); since also
\(d^*a_{\mathrm{coex}}=0\), the form \(a_{\mathrm{coex}}\) is harmonic, and being
coexact it must vanish.

The connected component of the Abelian gauge group acts by translation along the exact
summand:
\[
a\longmapsto a+d\phi,
\qquad
\phi\in \Omega^0(X;\mathfrak t).
\]
Thus quotienting by small gauge transformations removes the
\(d\Omega^0(X;\mathfrak t)\)-direction. If one parametrizes this direction by
\(\Omega^0(X;\mathfrak t)/\ker d\), then the map
\[
d:\Omega^0(X;\mathfrak t)/\ker d \longrightarrow d\Omega^0(X;\mathfrak t)
\]
has zeta-regularized Jacobian
\[
\det\nolimits'_{\zeta}(d^*d)^{1/2}
=
\det\nolimits'_{\zeta}(\Delta_{0,\mathfrak t})^{1/2}.
\]
This is the Faddeev--Popov factor. By Proposition~\ref{prop:moduli-torsor-toral},  the residual quotient by large gauge transformations identifies the harmonic
sector with the compact torus
\[
H^1(X;\mathfrak t)\big/ H^1(X;\Lambda)
\cong
\mathcal M_{X,p}(\mathbb T).
\]
Finally, the coexact sector contributes precisely the zeta-regularized oscillatory Gaussian
factor $\mathcal I_\zeta(B_{X,K}).$ Applying Proposition~\ref{prop:linear-quotient-factorization} to this Hodge splitting,
in the zeta-regularized Gaussian sense fixed above, we obtain
\[
\det\nolimits'_{\zeta}(\Delta_{0,\mathfrak t})^{1/2}\,
\mathcal I_\zeta(B_{X,K})\int_{\mathcal M_{X,p}(\mathbb T)} d\mu_{\mathrm{har}}\,,
\]
for the gauge-reduced integral over the fluctuation \(a\). Pulling out the constant factor $e^{2\pi i\,\operatorname{CS}_{X,P,K}(\Theta_p)}$
from the shift formula for the Chern--Simons action then yields
\eqref{eq:closed-gauge-reduction}.
\end{proof}

\end{document}
